\subsection{Bands in a fully lattice-ordered space}

DEFINITION 4. --- In a fully lattice-ordered space E, a linear subspace B of E is said to be a band if it satisfies the following conditions: 1) the relations \(x \in B\) , \(y \in E\) and \(|y| \leqslant |x|\) imply \(y \in B\) ; 2) for every nonempty subset X of B that is bounded above in E, the supremum sup X of X in E belongs to B.

Example. --- In the space \(R^{A}\) of real-valued functions defined on a set A, the set of functions that are zero at all the points of a subset M of A is a band.

\pandocbounded{\includegraphics[keepaspectratio]{images/part001_20a85035.jpg}}

Remark. --- In the space \(R^{A}\) , the subspace \(\mathcal{B}(A)\) of bounded real-valued functions on A satisfies condition 1) of Def. 4; moreover, for every subset X of \(\mathcal{B}(A)\) that is bounded above in \(\mathcal{B}(A)\) , the upper envelope of X belongs to \(\mathcal{B}(A)\) . However, if A is infinite, a subset of \(\mathcal{B}(A)\) may be bounded above in \(R^{A}\) without being bounded above in \(\mathcal{B}(A)\) , in which case \(\mathcal{B}(A)\) is not a band in \(R^{A}\) .

It follows at once from Def. 4 that if B is a band in E then, for every nonempty subset X of B that is bounded below in E, inf X belongs to B. Every band B in E, equipped with the ordered vector space structure induced by that of E, is a fully lattice-ordered space and, for every subset \(X \subset B\) that is bounded above in B, the supremum of X in B is identical with its supremum in E.

The intersection of any family of bands in a fully lattice-ordered space E is also a band. For every subset \(M \subset E\) , there exists a smallest band containing M (since E itself is a band); this band will be called the band generated by M.

The properties of bands in a fully lattice-ordered space rest on the following proposition:

PROPOSITION 5. --- Let E be a fully lattice-ordered space, A a non-empty subset of E consisting of elements \(\geqslant 0\) , such that: 1) \(A + A \subset A\) , and 2) the relations \(x \in A\) , \(0 \leqslant y \leqslant x\) imply \(y \in A\) . Let M be the set of suprema in E of those subsets of A that are bounded above in E. Under these conditions, every element \(x \geqslant 0\) of E may be written in the form \(y + z\) , where \(y \in M\) is the supremum of the elements \(v \in A\) such that \(v \leqslant x\) , and where z is an element \(\geqslant 0\) that is alien to every element of M.

At any rate \(y \leqslant x\) , so it all comes down to showing that z = x - y is alien to every element \(t \in A\) (No.~1), in other words that \(u = \inf(z, t)\) is zero. By hypothesis, \(u \in A\) and \(u \leqslant x - y\) , thus \(u + y \leqslant x\) ; for every \(v \in A\) such that \(v \leqslant x\) , by definition \(v \leqslant y\) , thus \(u + v \leqslant u + y \leqslant x\) ; since \(u + v \in A\) by hypothesis, also \(u + v \leqslant y\) by the definition of y; finally, since \(u + y\) is the supremum in E of the elements \(u + v\) such that \(v \in A\) and \(v \leqslant x\) , we have \(u + y \leqslant y\) , whence \(u \leqslant 0\) , which completes the proof.

THEOREM 1 (F. Riesz). --- Let A be a subset of a fully lattice-ordered space E. The set A' of elements that are alien to every element of A is a band; the band A'\,' of elements that are alien to every element of A' is identical to the band generated by A, and E is the ordered direct sum of the bands A' and A'\,'.

The properties of alien elements, reviewed in No.~1, and the definition of a band, show at once that A' is a band, hence so is A'\,'. By Proposition 5 and the definition of a band, every element \(x \geqslant 0\) of E may be written \(x = y + z\) , with \(y \in A'\) and \(z \in A''\) , y and z being \(\geqslant 0\) ; since every element of E is the difference of two elements \(\geqslant 0\) , we have \(E = A' + A''\) ; on the other hand, since 0 is the only element alien to itself, we have \(A' \cap A'' = \{0\}\) , which proves that E is the direct sum of \(A'\) and \(A''\) ; finally, since the components in \(A'\) and \(A''\) of an element \(\geqslant 0\) of E are \(\geqslant 0\) , E is the ordered direct sum of \(A'\) and \(A''\) (No.~4, Prop. 4).

It remains to show that \(A''\) is identical to the band B generated by A. Now, E is the direct sum of B and the band \(B'\) formed by the elements alien to all the elements of B; since \(A \subset B\) , we have \(B' \subset A'\) ; on the other hand \(B \subset A''\) and E is also the direct sum of \(A'\) and \(A''\) ; therefore necessarily \(B = A''\) , \(B' = A'\) .

Theorem 1 and Proposition 5 make it possible to give another definition of the band generated by a set of elements of E:

PROPOSITION 6. --- Let E be a fully lattice-ordered space, M a subset of E, and B the band generated by M. Let \(M_{1}\) be the set of elements \(\geqslant 0\) of E each of which is \(\leqslant\) some element of the form \(\sum |x_{i}|\) , where \(x_{i} \in M\) ;

let \(M_{2}\) be the set of suprema of those subsets of \(M_{1}\) that are bounded above; then the set \(M_{2}\) is identical with the set of elements \(\geqslant0\) of B.

Clearly \(M_{2} \subset B\) by the definition of a band; on the other hand, if \(B'\) is the band of elements that are alien to every element of \(M_{1}\) , Theorem 1 shows that E is the ordered direct sum of B and \(B'\) . But Proposition 5 shows that every element \(\geqslant 0\) of E is the sum of an element of \(M_{2}\) and an element of \(B'\) , whence the proposition.

COROLLARY. --- Let \(a\) be an element of a fully lattice-ordered space \(\mathbf{E}\) , \(\mathbf{B}_a\) the band generated by \(a\) , \(\mathbf{B}_a'\) the band of elements alien to \(a\) . For every element \(x \geqslant 0\) of \(\mathbf{E}\) , the component of \(x\) in \(\mathbf{B}_a\) (for the decomposition of \(\mathbf{E}\) as the ordered direct sum of \(\mathbf{B}_a\) and \(\mathbf{B}_a'\) ) is equal to \(\sup_{n \in \mathbf{N}} \left( \inf(n|a|, x) \right)\) .

This follows from Proposition 6, applied to \(\mathbf{M} = \{a\}\) , and Proposition 5.

Note that the bands generated by a and \(|a|\) are identical. If a and b are two elements of E that are alien to each other, and if A and B are the bands generated by a and b, respectively, then every element of A is alien to every element of B; for, b belongs to the band \(A'\) of elements alien to a, whence \(B \subset A'\) , and, by Theorem 1, every element of A is alien to every element of \(A'\) .

